\documentclass[10pt]{article}
\usepackage[utf8]{inputenc}
\usepackage[T1]{fontenc}
\usepackage{amsmath}
\usepackage{amsfonts}
\usepackage{amssymb}
\usepackage[version=4]{mhchem}
\usepackage{stmaryrd}
\usepackage{bbold}
\usepackage{hyperref}
\hypersetup{colorlinks=true, linkcolor=blue, filecolor=magenta, urlcolor=cyan,}
\urlstyle{same}

\begin{document}
\section*{Symbolic Evaluation of Polylogarithms at Special Points: Algebraic Structures, Functional Equations, and Experimental Mathematics}
\section*{Introduction to Polylogarithmic Evaluation and Algebraic Structures}
The polylogarithm function, denoted mathematically as $\operatorname{Li}_{n}(z)$, is defined by the infinite power series $\operatorname{Li}_{n}(z)=\sum_{k=1}^{\infty} \frac{z^{k}}{k^{n}} \quad$ for an arbitrary complex argument $z$ residing strictly within the open unit disk $|z|<1$. For arguments extending beyond this region, the function is defined via analytical continuation across the complex plane, typically manifesting a branch cut along the positive real axis from ${ }^{1}$ to positive infinity. While the function possesses a continuous domain over the complex numbers, its symbolic evaluation at specific "special points"-most notably rational numbers, algebraic numbers, roots of unity, and the roots of highly specific cubic and quartic polynomials-reveals profound structural symmetries that intersect multiple branches of modern mathematics. The existence of closed-form evaluations at these points is not merely a sequence of isolated analytical curiosities; rather, it reflects deep underlying connections to algebraic K-theory, the Bloch group, hyperbolic geometry, and the representation theory of affine Lie algebras. ${ }^{1}$\\
Historically, the evaluation of the dilogarithm (the case where $n=2$, denoted $\operatorname{Li}_{2}(z)$ ) at specific points was rigorously investigated by Leonhard Euler, John Landen, and William Spence, yielding a handful of foundational classical identities. ${ }^{1}$ Among these are the well-known values $\operatorname{Li}_{2}(1)=\frac{\pi^{2}}{6}, \operatorname{Li}_{2}(-1)=-\frac{\pi^{2}}{12}$, and the highly symmetric evaluation $\operatorname{Li}_{2}\left(\frac{1}{2}\right)=\frac{\pi^{2}}{12}-\frac{1}{2} \ln ^{2}(2) \quad$. ${ }^{4}$ For centuries, these evaluations were treated as rare artifacts of integration. However, the modern mathematical landscape, propelled by the advent of high-precision computational algorithms like the Partial Sum of Least Squares (PSLQ) algorithm, has facilitated the empirical discovery of highly complex, multi-term polylogarithmic identities that were previously entirely inaccessible to manual algebraic derivation. ${ }^{4}$ These numerically verified conjectures often demand sophisticated analytical machinery to prove, bridging the gap between computational experimental mathematics and abstract algebraic geometry. ${ }^{3}$ The analysis presented herein synthesizes advanced mathematical discussions concerning the symbolic evaluation of polylogarithms at special points across\\
scholarly forums, specifically extracting the underlying mathematical mechanics from the Math Stack Exchange and MathOverflow networks. By meticulously examining multi-term identities, parameterized functional equations, and the architecture of polylogarithmic ladders, this report extracts second and third-order insights into the mechanisms governing these special functions. The investigation comprehensively covers rational arguments, complex conjugate arguments, the architecture of polylogarithmic ladders rooted in cubic constants (such as the tribonacci constant), and the overarching algebraic frameworks that dictate their inevitable existence.

\section*{The Epistemology of Experimental Mathematics and Integer Relation Algorithms}
A defining characteristic of contemporary research into special function identities is the epistemological shift from deductive, top-down derivation to empirical, computational discovery followed by bottom-up analytical proof. Advanced integer relation algorithms, most notably PSLQ, allow mathematicians to numerically evaluate complex polylogarithmic expressions to tens of thousands, and sometimes hundreds of thousands, of decimal digits of precision. ${ }^{4}$ By selecting a carefully curated restricted pool of candidate terms-often constructed from Riemann zeta values, integer powers of logarithms of fundamental constants such as $\ln (2), \ln (3)$, and $\ln (\phi)$ (the golden ratio), alongside known transcendental numbers like ${ }^{\pi}$ and Catalan's constant ${ }^{G}$-researchers can computationally identify vanishing linear combinations with precise rational coefficients. ${ }^{4}$\\
This computational methodology is mathematically rigorous but presents severe logistical challenges. If the selected candidate pool is excessively broad, the PSLQ algorithm risks generating false positives due to numerical coincidence, or it may suffer from exponential time complexity, rendering the search computationally intractable. Consequently, mathematical intuition remains paramount in selecting the "siblings" of the target expression. These siblings are simpler analytical terms or structured variations (such as replacing occurrences of $\ln (2)$ with $\ln (3)$, or multiplying the candidate terms by algebraic integers like $\sqrt{3}$ ) that are hypothesized to share a structural homology in the Bloch group of the field under investigation. ${ }^{4}$\\
Once an identity is discovered and verified to an extraordinary precision-such as 30,000 to 100,000 decimal digits-it is elevated to the status of a formal conjecture. ${ }^{4}$ The rigorous proof of such conjectures typically requires reverse-engineering the arguments through the fundamental functional equations of the polylogarithm. This paradigm underscores a broader ripple effect in pure mathematics: computational brute force identifies the theoretical destination, while classical complex analysis and K-theory are subsequently employed to pave the rigorous, axiomatic path leading to it.

\section*{The Core Functional Equations of the Dilogarithm}
To understand the symbolic evaluation of the polylogarithm at special points, one must first deconstruct the primary functional equations that dictate its symmetry. The vast array of multi-term polylogarithmic identities is not a collection of disparate coincidences; rather, it represents the multifaceted algebraic manifestations of a singular underlying geometry. The functional relations of the dilogarithm are fundamentally derived from the five-term relation, originally formulated by Abel and subsequently generalized by Hill and others. ${ }^{1}$ The classical five-term identity for the dilogarithm can be written in various parameterizations, one of the most analytically useful being:

$$
\mathrm{Li}_{2}(z w)-\mathrm{Li}_{2}(w)-\mathrm{Li}_{2}(z)+\mathrm{Li}_{2}\left(\frac{w(1-z)}{1-w z}\right)+\mathrm{Li}_{2}\left(\frac{z(1-w)}{1-w z}\right)=-\ln \left(\frac{1-z}{1-w z}\right) \ln \left(\frac{1-w}{1-w z}\right)
$$

This equation serves as the algebraic cornerstone for evaluating multi-variable and single-variable identities. ${ }^{1}$ In advanced literature, it is often expressed utilizing the Rogers L-function, a normalized variant of the dilogarithm that absorbs the logarithmic residuals, maximizing the symmetry of the equations. ${ }^{1}$ Defined by Bytsko and others as $L(x)=\frac{6}{\pi^{2}}\left(\operatorname{Li}_{2}(x)+\frac{1}{2} \ln (x) \ln (1-x)\right)$, the Rogers L-function satisfies Abel's functional equation in a purely additive form without trailing logarithmic cross-terms. ${ }^{1}$ The formulations by Gordon, McIntosh, and Loxton offer a slight variation where\\
$L_{R}(x)=\mathrm{Li}_{2}(x)+\frac{1}{2} \ln (x) \ln (1-x) \quad$, which equals $\frac{\pi^{2}}{6} L(x) \quad$.\\
Under this framework, Euler's reflection relation reduces beautifully to $L(x)+L(1-x)=1$, and Abel's duplication formula translates to $\frac{1}{2} L\left(x^{2}\right)=L(x)-L(-x) \quad$. 1 It is this geometric rigidity that forces the polylogarithm to collapse into rational evaluations at specific nodes. Numbers $\theta \in(0,1)$ that satisfy equations of the form $\sum_{k=0}^{n} c_{k} L\left(\theta^{k}\right)=0$ for rational coefficients $c_{k}$ are formally classified as L-algebraic numbers. ${ }^{1}$ The study of these L-algebraic numbers dictates which algebraic fields can support polylogarithmic ladders.\\
Another indispensable functional equation is Landen's dilogarithm identity, which maps the domain inside the unit circle to the domain outside it, bridging the values of $z$ and $z /(z-1)$ :

$$
\operatorname{Li}_{2}(z)=-\operatorname{Li}_{2}\left(\frac{z}{z-1}\right)-\frac{1}{2} \ln ^{2}(1-z)
$$

This identity is crucial because it allows mathematicians to rotate complex arguments into purely imaginary or purely real domains. ${ }^{4}$ When combined with the inversion identity\\
$\operatorname{Li}_{2}(z)=-\operatorname{Li}_{2}(1 / z)-\frac{1}{2} \ln ^{2}(-z)-\frac{\pi^{2}}{6}$, these functional relationships form the complete toolkit required to dissect the empirical conjectures generated by PSLQ algorithms. ${ }^{8}$

\section*{Rational Arguments and High-Order Polylogarithms}
Relations connecting the values of the polylogarithm at rational points in the interval $(0,1)$ constitute a critical area of study, particularly as the order $n$ of the function $\mathrm{Li}_{n}(z)$ increases. For lower orders, such as the dilogarithm and trilogarithm, standard combinations of logarithms and Riemann zeta values are well-documented. Examples include vanishing linear combinations involving $\mathrm{Li}_{2}\left(\frac{1}{3}\right), \mathrm{Li}_{2}\left(\frac{1}{4}\right)$, and analogous arguments for $\mathrm{Li}_{3}$ and $\mathrm{Li}_{5}$. 4 For example, a straightforward application of the dilogarithmic ladder for the base ${ }^{3}$ yields the following identity involving rational points:

$$
6 \operatorname{Li}_{2}\left(\frac{1}{3}\right)-\operatorname{Li}_{2}\left(\frac{1}{9}\right)=\frac{\pi^{2}}{6}-\ln ^{2}(3)
$$

This relation is frequently utilized to explicitly zero out residual terms in multidimensional integral evaluations over the domain $[0, \infty)^{5}$ ? As the order of the polylogarithm increases, the complexity of the identities scales dramatically, generating massive polynomial coefficients. A striking example of this high-order phenomenon is an identity for the tetralogarithm ( ${ }^{\mathrm{Li}_{4}}$ ) discovered via PSLQ by Vladimir Reshetnikov:

$$
\begin{gathered}
720 \mathrm{Li}_{4}\left(\frac{1}{2}\right)-2160 \mathrm{Li}_{4}\left(\frac{1}{3}\right)+2160 \mathrm{Li}_{4}\left(\frac{2}{3}\right)+270 \mathrm{Li}_{4}\left(\frac{1}{4}\right)+540 \mathrm{Li}_{4}\left(\frac{3}{4}\right)+135 \mathrm{Li}_{4}\left(\frac{1}{9}\right) \\
=19 \pi^{4}+30\left(\pi^{2}-\ln ^{2}(3)\right)\left(10 \ln ^{2}(2)-12 \ln (2) \ln (3)+3 \ln ^{2}(3)\right)-30 \ln ^{2}(2)\left(19 \ln ^{2}(2)-24 \ln (2) \ln (3)+8 \ln ^{2}(3)\right)
\end{gathered}
$$

The proof of this tetralogarithmic identity requires navigating the structural properties of polylogarithmic ladders published by Leonard Lewin and leveraging the Fourier series of Bernoulli polynomials. By mapping the arguments to the inversion formula $\mathrm{Li}_{4}(z)=-\mathrm{Li}_{4}(1 / z)-\frac{16 \pi^{4}}{4!} B_{4}\left(\frac{\ln (z)}{2 \pi i}\right)$, the complex polynomial factors generated by the Bernoulli expansion precisely cancel down to recreate the massive integers on the right-hand side of the conjecture. ${ }^{4}$ Specifically, the leading $19 \pi^{4}$ term is directly derived from the evaluation of $-720 \frac{16 \pi^{4}}{4!} B_{4}\left(\frac{\ln (1 / 2)}{2 \pi i}\right)-135 \frac{16 \pi^{4}}{4!} B_{4}\left(\frac{\ln (1 / 9)}{2 \pi i}\right) \quad .4$ The most staggering demonstration of high-order rational evaluations, however, is found in the work of Don Zagier concerning the heptalogarithm ( ${ }^{\mathrm{Li}_{7}}$ ). Zagier constructed a highly intricate\\
function defined by the expansion of Bernoulli numbers $B_{j}$ :

$$
P_{7}\left(x_{n}\right)=\sum_{j=0}^{6} \frac{2^{j} B_{j}}{j!} \ln ^{j}\left|x_{n}\right| \operatorname{Li}_{7-j}\left(x_{n}\right)
$$

By utilizing a meticulously selected set of 29 rational arguments-ranging from simple fractions like $1 / 2$ and $-1 / 3$ to highly composite and unusual fractions like $-1 / 4374,-2 / 243$, and $32 / 81$-Zagier established a vanishing linear combination weighted by 29 specific, massive integer coefficients. ${ }^{4}$ The absolute largest of these coefficients reaches the magnitude of -151540388696 . The summation of these 29 weighted terms evaluates exactly to $\frac{-1020149599795}{96} \zeta(7) \quad{ }^{4}$ The existence of such an identity suggests that rational arguments of higher-order polylogarithms are governed by a hyper-dimensional lattice of algebraic relations. The explicit mapping of this lattice remains one of the preeminent open problems in the intersection of number theory and special functions.

\section*{Complex Arguments, Clausen Functions, and Hypergeometric Linkages}
The evaluation of polylogarithms at complex arguments, particularly those residing on the unit circle or representing specific roots of unity, relies heavily on Fourier series representations and Clausen functions. For an argument $z=e^{i \theta}$, the dilogarithm separates cleanly into real and imaginary components associated with the standard Clausen functions $\mathrm{Sl}_{2}(\theta)$ and $\mathrm{Cl}_{2}(\theta)$ :

$$
\mathrm{Li}_{2}\left(e^{i \theta}\right)=\mathrm{Sl}_{2}(\theta)+i \mathrm{Cl}_{2}(\theta)
$$

where $\mathrm{Sl}_{2}(\theta)=\sum_{k=1}^{\infty} \frac{\cos (k \theta)}{k^{2}}$ and $\mathrm{Cl}_{2}(\theta)=\sum_{k=1}^{\infty} \frac{\sin (k \theta)}{k^{2}}$. Utilizing the Fourier series of the Bernoulli polynomials, the real part $\mathrm{Sl}_{2}(\theta)$ simplifies algebraically to

$$
\frac{\pi^{2}}{6}-\frac{\pi \theta}{2}+\frac{\theta^{2}}{4} \quad \text { for } \theta \in[0,2 \pi)
$$

This analytical framework is instrumental in deriving closed forms for complex arguments where the modulus of $\frac{z}{z-1}$ equals unity. For instance, evaluating the real part of the complex dilogarithm $\operatorname{Li}_{2}\left(\frac{1}{2}+\frac{i}{6}\right)$ natively leverages Landen's dilogarithm identity. For $z=\frac{1}{2}+\frac{i}{6}$, the ratio $\frac{z}{z-1}$ algebraically reduces to $-\frac{4}{5}-\frac{3}{5} i$, which clearly lies on the unit circle in the complex plane. This point corresponds to a phase angle $\theta=\arctan (3 / 4)+\pi$.By\\
applying the Clausen function relationships, the real part is extracted without requiring direct integration, yielding exactly $\frac{7 \pi^{2}}{48}-\frac{1}{3} \arctan ^{2}(2)-\frac{1}{6} \arctan ^{2}(3)-\frac{1}{8} \ln ^{2}\left(\frac{18}{5}\right) \quad .4$ A similar methodology elegantly resolves conjectures involving algebraic irrationals multiplied by roots of unity. Consider the evaluation of $\operatorname{Li}_{2}\left(\sqrt{2-\sqrt{3}} \cdot e^{i \pi / 12}\right)$, which algebraically simplifies to $\operatorname{Li}_{2}\left(\frac{1}{2}+i\left(1-\frac{\sqrt{3}}{2}\right)\right)$. The parameter $\theta$ in this instance maps to $\frac{7 \pi}{6}$, necessitating the reduction of the Clausen function $\mathrm{Cl}_{2}\left(\frac{7 \pi}{6}\right)$. By grouping the sine terms of the Clausen series and utilizing the reflection, duplication, and triplication formulas of the polygamma function, $\mathrm{Cl}_{2}\left(\frac{7 \pi}{6}\right)$ is rigorously evaluated in terms of Catalan's constant $G$ and the trigamma function $\psi^{(1)}(1 / 3)$. ${ }^{4}$ When Gieseking's constant $V$ is introduced, recognizing that $\psi^{(1)}\left(\frac{1}{3}\right)=\frac{2}{3} \pi^{2}+2 \sqrt{3} V$, the evaluation achieves a profound symmetric form directly linked to the volume of specific hyperbolic manifolds. ${ }^{4}$

\section*{The Imaginary Golden Ratio and Complex Conjugate Symmetry}
The interplay between complex arguments and profound algebraic constants is further highlighted by evaluations involving the imaginary golden ratio, $i \phi$, where $\phi=\frac{1+\sqrt{5}}{2}$. Extensive numerical testing via PSLQ up to 20,000 decimal digits generated remarkable conjectures such as:

$$
\Re\left[1800 \operatorname{Li}_{3}(i \phi)-24 \operatorname{Li}_{3}\left(i \phi^{5}\right)\right]=100 \ln ^{3}(\phi)-47 \pi^{2} \ln (\phi)-150 \zeta(3)
$$

Analytical proofs of these conjectures require decomposing the even powers using the classical duplication identity $\mathrm{Li}_{3}\left(x^{2}\right)=4\left(\mathrm{Li}_{3}(x)+\mathrm{Li}_{3}(-x)\right) \quad .{ }^{4}$ This allows for the expression of $\Re \operatorname{Li}_{3}\left(i \phi^{k}\right)$ natively in terms of $\operatorname{Li}_{3}\left(-\phi^{-2 k}\right)$. By recursively applying the inversion identity $\operatorname{Li}_{3}(-x)-\operatorname{Li}_{3}(-1 / x)=-\zeta(2) \ln (x)-\frac{\ln ^{3}(x)}{6} \quad$ to the geometric powers of the golden ratio, the complex trilogarithms are translated entirely into real-valued trilogarithms of $\phi^{-2}$, which are known constants linked to the dilogarithmic ladders of Coxeter and Landen. ${ }^{4}$ This analytical mapping confirms that seemingly arbitrary complex evaluations of $\mathrm{Li}_{3}$ are mathematically tethered to the rigid algebraic topology of the golden ratio.\\
The evaluation of specialized complex identities often demands that conjugate symmetries be systematically untangled. When the real portion of a complex identity is sought, fundamental properties of real and imaginary projection (e.g., $\Re(i z)=-\Im(z)$ ) are applied to segregate\\
the dilogarithmic inputs. An excellent example resides in the evaluation of a special value dilogarithm identity: $c=\frac{1+i \sqrt{3}}{3} \operatorname{Li}_{2}\left(1-\frac{i \sqrt{3}}{3}\right)+\operatorname{Li}_{2}\left(\frac{3}{4}+\frac{i \sqrt{3}}{4}\right)+\ldots \quad$. The core of the evaluation relies on mapping the argument $z=\frac{3}{4}+i \frac{\sqrt{3}}{4}$ through Landen's identity. ${ }^{8}$ This functional mapping rotates the argument into a purely imaginary phase, specifically $-i \sqrt{3}$, shifting the analytical burden to evaluating $\operatorname{Li}_{2}(-i \sqrt{3})$. By subsequently invoking the mirror symmetry of the dilogarithm across the real axis, $\overline{\mathrm{Li}_{2}(z)}=\mathrm{Li}_{2}(\bar{z})$, the conjugate pairs are merged using the quadratic identity $\mathrm{Li}_{2}(z)+\mathrm{Li}_{2}(-z)=\frac{1}{2} \mathrm{Li}_{2}\left(z^{2}\right) \quad$. ${ }^{8}$ Consequently, an intricate web of complex arguments evaluates cleanly to a function of $\mathrm{Li}_{2}(-3)$, which via reciprocal inversion reduces to $\mathrm{Li}_{2}(-1 / 3)$ plus known logarithmic constants. ${ }^{8}$ The apparent complexity of the imaginary plane collapses entirely under the weight of Landen's identity, translating geometric rotation into elementary algebraic reciprocation.

\section*{Hypergeometric Linkages to Complex Trilogarithms}
When evaluating the imaginary part of $\operatorname{Li}_{3}\left(\frac{1+i}{2}\right)$, the quest for a closed form reveals a strict limitation in purely elementary functions, necessitating the introduction of generalized hypergeometric functions. ${ }^{11}$ Polylogarithms are intimately related to hypergeometric series via the fundamental relation $\operatorname{Li}_{n}(z)=z_{n+1} F_{n}(1,1, \ldots, 1 ; 2,2, \ldots, 2 ; z) \quad$.\\
Through the manipulation of the inverse tangent integral of order 3 , defined as $\mathrm{Ti}_{3}(x)=\Im\left[\mathrm{Li}_{3}(i x)\right]$, the imaginary component of $\mathrm{Li}_{3}\left(\frac{1+i}{2}\right)$ can be definitively expressed as:

$$
\Im\left[\operatorname{Li}_{3}\left(\frac{1+i}{2}\right)\right]=\frac{\pi^{3}}{128}+\frac{\pi}{32} \ln ^{2}(2)+\frac{1}{4}{ }_{4} F_{3}\left(\frac{1}{2}, \frac{1}{2}, 1,1 ; \frac{3}{2}, \frac{3}{2}, \left.\frac{3}{2} \right\rvert\, 1\right)
$$

This remarkable reduction establishes a critical analytical linkage: the algebraic structure of the trilogarithm evaluated exactly at the midpoint of the complex unit square is inextricably bound to the evaluation of the ${ }_{4} F_{3}$ hypergeometric function at unity. ${ }^{11}$ Furthermore, because the sum $\Im\left[\operatorname{Li}_{3}\left(\frac{1+i}{2}\right)\right]+\Im\left[\operatorname{Li}_{3}(1+i)\right] \quad$ evaluates algebraically to exactly $\frac{7 \pi^{3}}{128}+\frac{3 \pi}{32} \ln ^{2}(2)$ , solving the functional evaluation for one point immediately resolves the other. This demonstrates the symmetric rigidity of the trilogarithm in the complex plane, confirming that hypergeometric representations are not merely approximations but exact boundary conditions of the function. ${ }^{11}$

\section*{The Architecture of Polylogarithmic Ladders}
A "polylogarithmic ladder" represents one of the most structurally significant and visually striking phenomena in the study of special functions. A ladder of weight $n$ and index $N$ evaluated at a $\mathbb{Q}_{\text {-algebraic argument }} u$ is a specific linear combination taking the form:

$$
\sum_{r=0}^{N-1} A_{r} \frac{\operatorname{Li}_{n}\left(u^{r}\right)}{r^{n-1}}
$$

where the coefficients $A_{r}$ are entirely rational, and the right-hand side of the equation reduces completely to a linear combination of products of simple logarithms and lower-weight Riemann zeta values. ${ }^{1}$ The existence of these ladders is profoundly non-trivial. They only occur when the geometric argument $u$ is a root of specific algebraic equations (most commonly cubic or quartic polynomials) that possess exceptional properties in the Bloch group of the corresponding algebraic number field.

\section*{Cubic Constants: Tribonacci, Plastic, and Supergolden Ratios}
The dilogarithmic and trilogarithmic ladders associated with specific cubic algebraic constants showcase extraordinary mathematical symmetry. The most prominent among these constants is the tribonacci constant $\tau$, the unique real root of the cubic equation $x^{3}-x^{2}-x-1=0$, which possesses a numerical value of approximately 1.839286 A conjectured dilogarithm identity containing the tribonacci constant posits that:

$$
\operatorname{Li}_{2}\left(\frac{1}{\tau^{3}}\right)+\operatorname{Li}_{2}\left(\tau^{2}\right)+\operatorname{Li}_{2}\left(\frac{\tau}{\tau+1}\right)=\frac{1}{2} \ln (1-\tau) \ln (\tau)-\cdots-\frac{7 \pi^{2}}{12}
$$

The analytical proof of this specific identity leverages the concept of functional equivalence, denoted by \~{} , which implies that the two sides of an equation differ exclusively by elementary logarithmic functions. ${ }^{7}$ Applying standard inversion and reflection identities transforms the distinct geometric terms $\mathrm{Li}_{2}\left(\frac{\tau}{\tau+1}\right)$ and $\mathrm{Li}_{2}\left(\tau^{2}\right)$ into expressions involving strictly $\tau^{-1}$ and $\tau^{-2}$. Let $\sigma=\tau^{-1}$ (the inverse tribonacci constant, satisfying the symmetric equation $\sigma^{3}+\sigma^{2}+\sigma=1$ ). Under this substitution, the identity collapses into a pure dilogarithmic ladder of the form $\operatorname{Li}_{2}\left(\sigma^{3}\right)-\frac{1}{2} \operatorname{Li}_{2}\left(\sigma^{2}\right)-\operatorname{Li}_{2}(\sigma) \simeq 0$. This specific ladder is a direct, linear combination of the structural ladders mapped in Leonard Lewin's foundational text, Structural Properties of Polylogarithms. ${ }^{7}$ Another analytical proof utilizes Hill's 5-term\\
dilogarithm identity with variables mapped directly to $w=u^{3}$ and $z=-u^{-1}$, proving the relation strictly via classical substitutions. ${ }^{7}$\\
The phenomenon extends deeply into the third mathematical order. Trilogarithm ladders exist for a highly specific, curated quartet of cubic constants: the tribonacci constant ( $\left.x^{3}=x^{2}+x+1\right)$, the supergolden ratio ( $y^{3}=y^{2}+1$ ), the plastic ratio ( $z^{3}=z+1$ ), and the supersilver ratio ( $w^{3}=2 w^{2}+1$ ). ${ }^{12}$ The empirical discovery of these trilogarithmic ladders prompts the fundamental question of whether their existence is a trivial outcome of polynomial manipulation or indicative of deeper algebraic geometry. Analytical proofs demonstrate that these ladders are generated deterministically by two fundamental functional equations of the trilogarithm:

\begin{enumerate}
  \item $\mathrm{Li}_{3}(X)+\mathrm{Li}_{3}(1-X)+\mathrm{Li}_{3}\left(1-X^{-1}\right)=\zeta(3)+\frac{1}{6} \ln ^{3}(X)+\frac{\pi^{2}}{6} \ln (X)-\frac{1}{2} \ln ^{2}(X) \ln (1-X)$
  \item $\mathrm{Li}_{3}(X)+\mathrm{Li}_{3}(-X)=\frac{1}{4} \mathrm{Li}_{3}\left(X^{2}\right)$
\end{enumerate}

By substituting the inverse of the supergolden ratio $\left(X=y^{-1}\right)$ into the first functional identity, and utilizing the defining algebraic relation $y^{3}=y^{2}+1$ to simplify complex fractional terms like $1-y^{-1}$ directly into $y^{-3}$, the trilogarithms elegantly fold into straightforward powers of the base constant. ${ }^{12}$ The subsequent application of the duplication formula (Identity 2 ) to eliminate the negative arguments yields a perfectly balanced trilogarithm ladder that evaluates to unity when properly normalized against its logarithmic residuals and $\zeta(3){ }_{.}^{12}$\\
A parallel mechanism applies rigorously to the plastic ratio ( $X=z^{-2}$ ). Substituting $X=z^{-2}$ into Identity 1 forces the term $1-X$ to evaluate to $z^{-3}$ and $1-X^{-1}$ to evaluate to $z^{-1}$ due to the polynomial constraint $z\left(z^{2}-1\right)=1$ derived from $z^{3}=z+1 \quad .^{12}$ This proves conclusively that these ladders are not random; they are deterministic mathematical outcomes of the algebraic roots satisfying the specific constraint that $1-X^{k}=X^{j}$, allowing the functional equations to loop recursively over themselves. ${ }^{12}$

\section*{Advanced Algebraic Identities and Parameterized Functional Equations}
As analytical capabilities have expanded, the search for special values has moved beyond static roots into parameterized continuous identities, confirming that continuous functional relations of the dilogarithm are ultimately shadows of the primary five-term relation.\\
An example of this continuous symmetry is highlighted by a massive multi-term equation\\
involving the parameter $m$. The identity features expressions of immense complexity, such as:

$$
\operatorname{Li}_{2}\left(\frac{2 m}{m^{2}+m-\sqrt{((m-3) m+1)\left(m^{2}+m+1\right)}-1}\right)
$$

combined with reciprocal terms and trailing logarithmic multipliers equal to zero for a real number $m$ in a neighborhood of $1^{6}$. While visually intimidating, the structural analysis of this parameterized identity reveals simple underlying algebraic constraints. By identifying the complex arguments of the four distinct dilogarithms as continuous variables $a_{1}, a_{2}, a_{3}, a_{4}$, researchers established the elementary relations $\left(1-a_{1}\right) a_{4}=1$ and $a_{2} a_{3}=-1+a_{3} \quad$. ${ }^{6}$ Because of these relationships, applying Euler's reflection formula followed by an inversion formula natively reduces the entire sprawling dilogarithm expression into a trivial zero, without ever needing to rely on calculus or integration. ${ }^{6}$

\section*{The $\pi / 9$ Identities and Kirillov's Algebraic Proof}
A historical benchmark in the rigorous study of dilogarithmic evaluation is the set of identities involving the secants of multiples of $\pi / 9$. Specifically, the highly celebrated identity:

$$
\mathrm{Li}_{2}(-a)+\mathrm{Li}_{2}\left(a^{2}\right)-\frac{1}{3} \mathrm{Li}_{2}\left(-a^{3}\right)=-\frac{\pi^{2}}{54}
$$

where the geometric argument is $a=2 \cos (4 \pi / 9)$, forms one-third of a triad of identities originally explored by Loxton and Lewin. ${ }^{1}$ These specific bases are intrinsically tied to the roots of the cubic algebraic equation $x^{3}+3 x^{2}-1=0$.\\
For decades, the third identity in this triad existed solely as a numerical fact verified by computers, lacking any form of analytical derivation. It was not until 1990 that H. Gangl proved it utilizing methods based on Bloch groups. Subsequently, in 1995, Anatol N. Kirillov published an exhaustive, groundbreaking algebraic proof relying exclusively on the Rogers dilogarithm five-term relation. ${ }^{1}$ Kirillov's work established definitively that the seemingly arbitrary angles of $\pi / 9$ and $2 \pi / 9$ are geometric necessities dictated by the algebraic properties of the roots of unity $e^{2 \pi i / 9}$ when constrained by the algebraic geometry of the five-term relation. ${ }^{1}$ More recent, modernized analytical proofs bypass Kirillov's exceptionally dense algebraic maneuvers by employing a two-part trigonometric and polynomial mapping approach. By substituting roots of the specific polynomial $u^{3}-3 u^{2}+2 u+1 / 3=0 \quad$ into Landen's\\
identity, and equating the resulting radical expressions to the complex boundary $i \sqrt{3}$, the complex logarithm arguments are forced computationally to map precisely to the minimal polynomial for $\cos (\pi / 9)$. Consequently, the complex part inside the logarithm simplifies naturally to $\pm \sin (\pi / 9)$, confirming the $-\frac{\pi^{2}}{54}$ evaluation with startling analytical efficiency. ${ }^{1}$

\section*{Beta Integral Series Evaluation (Hakimoglu-Brown Technique)}
An emerging modern technique for resolving outstanding complex dilogarithm conjectures involves beta integral-based series evaluations. Formulated heavily in the contemporary work of Hakimoglu-Brown, this methodology relies explicitly on evaluating "sextic integrals" (integrals featuring a precise sixfold parameter combination) of the form:

$$
\int_{0}^{h} \frac{(1-y)^{c} y^{d}}{1+g^{2}(1-y)^{a} y^{b}} d g
$$

By mapping these integrals mathematically to generalized hypergeometric series ${ }_{4} F_{3}$ and employing the dominated convergence theorem, mathematicians can extract shifting relations directly without resorting to the cumbersome double integration techniques utilized historically by scholars like Batir. ${ }^{1}$\\
This methodology has successfully yielded single-term evaluations of exceptional computational complexity, such as an 8th-degree geometric identity involving $y \approx 1.3945-0.6970 i$ (a verified root of $y^{8}-3 y^{7}+\cdots+1=0$ ). ${ }^{1}$ Furthermore, it has permitted the mathematical construction of completely novel three-term ladder identities situated within complex algebraic number fields like $\mathbb{Q}(\sqrt{-7})$, vastly expanding the known taxonomy of L-algebraic numbers. ${ }^{1}$

\section*{Extracted Repository of Relevant Mathematical}
\section*{Queries}
To contextualize the scope of these symbolic evaluations, a comprehensive review of the Math Stack Exchange and MathOverflow networks was conducted. The following curated repository catalogs the specific analytical queries that fulfill the strict criteria of this investigation. Each entry represents a query that explicitly contains a polylogarithmic identity evaluated at a special point within its question body, and possesses at least one mathematically rigorous analytical solution. Questions lacking explicit mathematical identities in the prompt or those devoid of meaningful analytical answers have been rigorously excluded.

\begin{center}
\begin{tabular}{|l|l|}
\hline
Source URL & Question Title \\
\hline
\end{tabular}
\end{center}

\begin{center}
\begin{tabular}{|l|l|}
\hline
https://math.stackexchange.com/questions/ 1424600/... & Conjecture Re(Li\_2(1/2+i/6)) = 7\textbackslash pi\^{}2/48-1/3 $\arctan ^{\wedge} 22-1 / 6 \arctan ^{\wedge} 23-1 / 8 \ln ^{\wedge} 2(18 / 5)$ \\
\hline
https://math.stackexchange.com/questions/ 1430169/... & Conjectured closed form for Li\_2(|sqrt\{2-|sqrt\{3\}\}|cdot e\^{}\{i|pi/12\}) \\
\hline
https://math.stackexchange.com/questions/ 1373123/... & A conjectured identity for tetralogarithms Li\_4 \\
\hline
https://math.stackexchange.com/questions/ 945972/... & Relations connecting values of the polylogarithm Li\_n at rational points \\
\hline
https://math.stackexchange.com/questions/ 1579355/... & Trilogarithm Li\_3(z) and the imaginary golden ratio ilphi \\
\hline
https://math.stackexchange.com/questions/ 5122160/... & Is there a closed form for Re(Li\_2(3i)) Nielsen polylogarithm \\
\hline
https://math.stackexchange.com/questions/ 4284361/... & Closed form evaluation of a trigonometric integral in terms of polylogarithms \\
\hline
https://math.stackexchange.com/questions/ 918680/... & Closed Form for the Imaginary Part of Li\_3((1+i)/2) \\
\hline
https://math.stackexchange.com/questions/ 812194/... & Dilogarithm in closed form \\
\hline
https://math.stackexchange.com/questions/ 980179/... & Closed-form of a special value dilogarithm identity \\
\hline
https://math.stackexchange.com/questions/ 1411312/... & Closed form of a series dilogarithm \\
\hline
https://math.stackexchange.com/questions/ 5125161/... & Verify a polylogarithmic evaluation of a 5D integral over A surface-level evaluation might attempt to solve this system via integral calculus, incorrectly expressing it directly as the indefinite integral \\
\hline
\end{tabular}
\end{center}

\begin{center}
\begin{tabular}{|l|l|}
\hline
 & $\int_{\sqrt{2}-1}^{1} \frac{\ln (t)}{1-t^{2}} d$ and attempting to solve it via integration by parts. ${ }^{5}$ However, the calculus approach fundamentally obscures the underlying algebraic symmetry dictating the roots. \\
\hline
\end{tabular}
\end{center}

By initiating the analytical proof via Abel's identity explicitly applied to the algebraic root $z=1-\sqrt{2}$, and subsequently threading the resultant expression through Landen's identity, the evaluation bypasses the necessity for integration entirely. ${ }^{5}$ It confirms mathematically that the identity reduces exactly to the closed form\\
$\frac{\pi^{2}}{8}-\frac{1}{8} \ln ^{2}(2)-\frac{1}{2} \ln ^{2}(1+\sqrt{2}) \quad .^{5}$ The mathematical ability to generalize this specific root relation to any arbitrary complex number $z_{\text {where }}|1-z|<1$ and $z \notin(-\infty, 0]$ confirms a critical tertiary insight: special values are not isolated phenomena. They are merely the discrete intersections where a continuous functional equation governing the complex plane generates highly stable rational or quadratic algebraic outputs. ${ }^{5}$\\
To evaluate the closed form of the golden ratio $\varphi_{\text {and its algebraic conjugate }} \Phi=\varphi-1$, analytical methods rely heavily on the pentagonal identity. ${ }^{14}$ Known to mathematicians as the most functionally significant dilogarithm identity, it takes the form:

$$
\operatorname{Li}_{2}\left(\frac{x}{1-y}\right)+\operatorname{Li}_{2}\left(\frac{y}{1-x}\right)-\operatorname{Li}_{2}(x)-\operatorname{Li}_{2}(y)-\operatorname{Li}_{2}\left(\frac{x y}{(1-x)(1-y)}\right)=\ln (1-x) \ln (1-y)
$$

By defining $\alpha=\frac{\sqrt{5}-1}{2}$ and substituting $x=y=1-\alpha$ into the pentagonal identity, the relation collapses into a linear combination of $\mathrm{Li}_{2}(\alpha)$ and $\mathrm{Li}_{2}(1-\alpha)$. ${ }^{14}$ By generating a precise system of three linear equations utilizing the reflection formula and the duplication identity, the precise values are algebraically isolated: $\operatorname{Li}_{2}\left(\frac{\sqrt{5}-1}{2}\right)=\frac{\pi^{2}}{10}-\ln ^{2}\left(\frac{\sqrt{5}-1}{2}\right)$, confirming the exact boundary parameters of the golden ratio's topology. ${ }^{14}$

\section*{Conclusion}
The symbolic evaluation of polylogarithms at special points sits at the exact intersection of algebraic number theory, experimental computational modeling, and classical complex analysis. The analytical data indicates a clear, unified architectural hierarchy: at the foundational mathematical level, the five-term relation and its associated Bloch group homologies govern all continuous functional behavior. When these continuous functional relations are restricted analytically to specific algebraic numbers-specifically roots of unity and the roots of select cubic and quartic polynomials like the tribonacci and plastic constants-they produce the discrete, mathematically rigid structures known to researchers as polylogarithmic ladders.

The integration of high-precision computational algorithms like PSLQ has fundamentally altered the trajectory of this mathematical field. It has successfully enabled the discovery of multi-term vanishing combinations and complex rational point evaluations (such as Don Zagier's 29-term heptalogarithm sequence) that possess a combinatorial complexity far exceeding the reach of unassisted human intuition. However, as demonstrated extensively by the subsequent application of hypergeometric series transformations, Clausen function sine expansions, and sextic beta integral evaluations, the requirement for absolute analytical rigor remains absolute. The identities highlighted throughout this research do not exist as mere numerical artifacts of computational power; they are explicit, provable boundary conditions of hyperbolic manifolds and the geometric representations of affine Lie algebras. Understanding these special points is, therefore, not merely an exercise in evaluating a power series, but an active endeavor in mapping the deepest algebraic symmetries inherent in modern theoretical mathematics.

\section*{Works cited}
\begin{enumerate}
  \item nt.number theory - Proving identity involving dilogarithms and $\$ 1 \mathrm{pi} / 9$..., accessed June 1, 2026, \href{https://mathoverflow.net/questions/492746/proving-identity-involving-dilogarithm}{https://mathoverflow.net/questions/492746/proving-identity-involving-dilogarithm} s-and-pi-9
  \item Polylogarithm sheaves - ag.algebraic geometry - MathOverflow, accessed June 1, 2026, \href{https://mathoverflow.net/questions/242920/polylogarithm-sheaves}{https://mathoverflow.net/questions/242920/polylogarithm-sheaves}
  \item What is special about polylogarithms that leads to so many interesting identities and applications? - MathOverflow, accessed June 1, 2026, \href{https://mathoverflow.net/questions/25428/what-is-special-about-polylogarithms}{https://mathoverflow.net/questions/25428/what-is-special-about-polylogarithms} -that-leads-to-so-many-interesting-identitie
  \item mathse-polylog-rendered-posts\_\_71a4599d00e3.pdf
  \item calculus - Proof of a dilogarithm identity - Mathematics Stack Exchange, accessed June 1, 2026, \href{https://math.stackexchange.com/questions/1959503/proof-of-a-dilogarithm-iden}{https://math.stackexchange.com/questions/1959503/proof-of-a-dilogarithm-iden} tity
  \item special functions - a dilogarithm identity: known or new ..., accessed June 1, 2026, \href{https://mathoverflow.net/questions/85986/a-dilogarithm-identity-known-or-new}{https://mathoverflow.net/questions/85986/a-dilogarithm-identity-known-or-new}
  \item calculus - Dilogarithm identity containing the tribonacci constant ..., accessed June 1, 2026, \href{https://math.stackexchange.com/questions/1436493/dilogarithm-identity-containi}{https://math.stackexchange.com/questions/1436493/dilogarithm-identity-containi} ng-the-tribonacci-constant
  \item calculus - Closed-form of a special value dilogarithm identity ..., accessed June 1, 2026, \href{https://math.stackexchange.com/questions/980179/closed-form-of-a-special-val}{https://math.stackexchange.com/questions/980179/closed-form-of-a-special-val} ue-dilogarithm-identity
  \item Verify a polylogarithmic evaluation of a 5D integral over $[0, \infty) 5$ - Math Stack Exchange, accessed June 1, 2026, \href{https://math.stackexchange.com/questions/5125161/verify-a-polylogarithmic-eval}{https://math.stackexchange.com/questions/5125161/verify-a-polylogarithmic-eval} uation-of-a-5d-integral-over-0-infty5
  \item A conjectured identity for tetralogarithms \$loperatorname\{Li\}\_4 - Math Stack
\end{enumerate}

Exchange, accessed June 1, 2026, \href{https://math.stackexchange.com/questions/1373123/a-conjectured-identity-for-t}{https://math.stackexchange.com/questions/1373123/a-conjectured-identity-for-t} etralogarithms-operatornameli-4\\
11. integration - Closed Form for the Imaginary Part of \$\textbackslash text\{Li\}\_3\textbackslash Big ..., accessed June 1, 2026, \href{https://math.stackexchange.com/questions/918680/closed-form-for-the-imagina}{https://math.stackexchange.com/questions/918680/closed-form-for-the-imagina} ry-part-of-textli-3-big-frac1i2-big\\
12. special functions - Why does the tribonacci constant have a ..., accessed June 1, 2026, \href{https://math.stackexchange.com/questions/1064100/why-does-the-tribonacci-c}{https://math.stackexchange.com/questions/1064100/why-does-the-tribonacci-c} onstant-have-a-trilogarithm-ladder\\
13. sequences and series - Polylogarithm ladders for the tribonacci and ..., accessed June 1, 2026, \href{https://math.stackexchange.com/questions/481682/polylogarithm-ladders-for-th}{https://math.stackexchange.com/questions/481682/polylogarithm-ladders-for-th} e-tribonacci-and-n-nacci-constants\\
14. real analysis - Mathematics Stack Exchange - Math Stack Exchange, accessed June 1, 2026, \href{https://math.stackexchange.com/questions/463682/closed-form-of-operatornam}{https://math.stackexchange.com/questions/463682/closed-form-of-operatornam} eli-2-varphi-and-operatornameli-2-varphi-1


\end{document}